\documentclass[11pt,a4paper]{article}
\usepackage[utf8]{inputenc}
\usepackage[spanish]{babel}
\usepackage[T1]{fontenc}
\usepackage{geometry}
\geometry{margin=2cm}
\usepackage{xcolor}
\usepackage{listings}
\usepackage{tcolorbox}
\tcbuselibrary{most}
\usepackage{booktabs}
\usepackage{fancyhdr}
\usepackage{titlesec}
\usepackage{multicol}

\definecolor{fapigreen}{RGB}{0,150,100}
\definecolor{codebg}{RGB}{245,245,245}
\definecolor{alertred}{RGB}{220,50,50}
\definecolor{tipgreen}{RGB}{39,174,96}
\definecolor{warnambel}{RGB}{230,126,34}
\definecolor{darkgray}{RGB}{50,50,50}

\lstdefinestyle{python}{
  language=Python,
  backgroundcolor=\color{codebg},
  basicstyle=\ttfamily\small,
  keywordstyle=\color{fapigreen}\bfseries,
  stringstyle=\color{teal},
  commentstyle=\color{gray}\itshape,
  numbers=left, numberstyle=\tiny\color{gray},
  frame=single, framerule=0.4pt, rulecolor=\color{fapigreen!40},
  breaklines=true, tabsize=2, showstringspaces=false,
  morekeywords={Optional,List,Union,Dict,Any,BaseModel,Field,validator,
                Depends,HTTPException,status,BackgroundTasks,Request,Response}
}

\newtcolorbox{tipbox}{colback=tipgreen!10,colframe=tipgreen,title=\textbf{Tip},fonttitle=\small\bfseries}
\newtcolorbox{alertbox}{colback=alertred!10,colframe=alertred,title=\textbf{Ojo},fonttitle=\small\bfseries}
\newtcolorbox{ejerciciobox}[1]{colback=warnambel!8,colframe=warnambel,title=\textbf{Ejercicio #1},fonttitle=\small\bfseries}
\newtcolorbox{solucionbox}[1]{colback=tipgreen!8,colframe=tipgreen!60,title=\textbf{Solución #1},fonttitle=\small\bfseries}

\pagestyle{fancy}
\fancyhf{}
\lhead{\textcolor{fapigreen}{\textbf{FastAPI en Profundidad}}}
\rhead{\textcolor{gray}{Preparación Técnica — Toku}}
\cfoot{\thepage}

\titleformat{\section}{\large\bfseries\color{fapigreen}}{}{0em}{}[\titlerule]
\titleformat{\subsection}{\normalsize\bfseries\color{darkgray}}{}{0em}{}

\begin{document}

\begin{titlepage}
  \centering
  \vspace*{2cm}
  {\Huge\bfseries\color{fapigreen} FastAPI en Profundidad\par}
  \vspace{0.5cm}
  {\large\color{gray} Más allá de los tutoriales — lo que Toku usa en producción\par}
  \vspace{0.3cm}
  \textcolor{fapigreen!40}{\rule{8cm}{2pt}}
  \vspace{1.5cm}

  \begin{tcolorbox}[colback=fapigreen!5,colframe=fapigreen,width=11cm]
    \centering
    Ya sabes Python. Ya hiciste algo con FastAPI.\\[4pt]
    Este documento te lleva al nivel de producción.
  \end{tcolorbox}
  \vspace{1.5cm}
  \begin{tabular}{ll}
    \textcolor{fapigreen}{\textbullet} & Pydantic v2 para validación\\
    \textcolor{fapigreen}{\textbullet} & Dependency Injection (el corazón de FastAPI)\\
    \textcolor{fapigreen}{\textbullet} & Async y manejo de conexiones DB\\
    \textcolor{fapigreen}{\textbullet} & Middleware y manejo de errores\\
    \textcolor{fapigreen}{\textbullet} & Estructura de proyecto real\\
    \textcolor{warnambel}{\textbullet} & Ejercicios de diseño de API\\
  \end{tabular}
  \vfill
  {\small\color{gray} Marzo 2026 — Estudio Personal\par}
\end{titlepage}

\tableofcontents
\newpage

% ─── SECCIÓN 1: PYDANTIC V2 ──────────────────────────────────────────────────
\section{Pydantic v2 — Validación y Serialización}

\begin{lstlisting}[style=python]
from pydantic import BaseModel, Field, field_validator, model_validator
from typing import Optional
from datetime import date
from enum import Enum

class Moneda(str, Enum):
    CLP = "CLP"
    MXN = "MXN"
    BRL = "BRL"

class InvoiceCreate(BaseModel):
    customer_id: str
    monto: int = Field(gt=0, description="Monto en centavos, debe ser positivo")
    moneda: Moneda = Moneda.CLP
    fecha_vencimiento: date
    external_id: Optional[str] = None
    descripcion: str = Field(default="", max_length=500)

    @field_validator("external_id")
    @classmethod
    def external_id_no_espacios(cls, v):
        if v and " " in v:
            raise ValueError("external_id no puede tener espacios")
        return v

    @model_validator(mode="after")
    def vencimiento_futuro(self):
        from datetime import date as D
        if self.fecha_vencimiento <= D.today():
            raise ValueError("La fecha de vencimiento debe ser futura")
        return self

class InvoiceResponse(BaseModel):
    id: str
    customer_id: str
    monto: int
    moneda: Moneda
    estado: str
    fecha_vencimiento: date
    created_at: str

    model_config = {"from_attributes": True}  # para leer desde ORM/asyncpg
\end{lstlisting}

\begin{tipbox}
\texttt{model\_config = \{"from\_attributes": True\}} reemplaza al viejo \texttt{orm\_mode = True} de Pydantic v1.
Es necesario para crear modelos desde objetos de base de datos (asyncpg Records).
\end{tipbox}

% ─── SECCIÓN 2: DEPENDENCY INJECTION ─────────────────────────────────────────
\section{Dependency Injection — El Corazón de FastAPI}

\begin{lstlisting}[style=python]
# dependencies.py
from fastapi import Depends, HTTPException, status
from fastapi.security import HTTPBearer, HTTPAuthorizationCredentials
import asyncpg

# Pool de conexiones a la BD
async def get_db() -> asyncpg.Connection:
    """Dependency: proporciona una conexion de BD por request"""
    async with db_pool.acquire() as conn:
        yield conn  # yield = context manager; libera la conexion al terminar

# Autenticacion
security = HTTPBearer()

async def get_current_org(
    credentials: HTTPAuthorizationCredentials = Depends(security),
    db: asyncpg.Connection = Depends(get_db),
) -> dict:
    """Dependency: verifica el API key y retorna la organizacion"""
    api_key = credentials.credentials

    org = await db.fetchrow(
        "SELECT id, nombre, pais FROM organizations WHERE api_key = $1",
        api_key
    )
    if not org:
        raise HTTPException(
            status_code=status.HTTP_401_UNAUTHORIZED,
            detail="API Key invalida"
        )
    return dict(org)

# En el router:
@router.post("/invoices", response_model=InvoiceResponse)
async def crear_invoice(
    payload: InvoiceCreate,
    org: dict = Depends(get_current_org),  # autenticacion automatica
    db: asyncpg.Connection = Depends(get_db),  # BD automatica
):
    # org ya fue validada por get_current_org
    invoice = await db.fetchrow(
        """INSERT INTO invoices (customer_id, monto, moneda, fecha_vencimiento, external_id)
           VALUES ($1, $2, $3, $4, $5)
           RETURNING *""",
        payload.customer_id, payload.monto, payload.moneda.value,
        payload.fecha_vencimiento, payload.external_id
    )
    return InvoiceResponse(**dict(invoice))
\end{lstlisting}

% ─── SECCIÓN 3: MANEJO DE ERRORES ────────────────────────────────────────────
\section{Manejo de Errores Consistente}

\begin{lstlisting}[style=python]
# errors.py
from fastapi import FastAPI, Request
from fastapi.responses import JSONResponse
from pydantic import ValidationError
import asyncpg

app = FastAPI()

class TokuAPIError(Exception):
    def __init__(self, mensaje: str, codigo: str, status_code: int = 400):
        self.mensaje = mensaje
        self.codigo = codigo
        self.status_code = status_code

@app.exception_handler(TokuAPIError)
async def manejar_toku_error(request: Request, exc: TokuAPIError):
    return JSONResponse(
        status_code=exc.status_code,
        content={"error": exc.codigo, "mensaje": exc.mensaje}
    )

@app.exception_handler(asyncpg.UniqueViolationError)
async def manejar_duplicado(request: Request, exc: asyncpg.UniqueViolationError):
    return JSONResponse(
        status_code=409,
        content={"error": "DUPLICATE", "mensaje": "Ya existe un registro con ese ID"}
    )

@app.exception_handler(ValidationError)
async def manejar_validacion(request: Request, exc: ValidationError):
    return JSONResponse(
        status_code=422,
        content={"error": "VALIDATION_ERROR", "detalle": exc.errors()}
    )

# Uso en ruta:
async def obtener_cliente(id: str, db=Depends(get_db)):
    cliente = await db.fetchrow("SELECT * FROM customers WHERE id = $1", id)
    if not cliente:
        raise TokuAPIError(
            mensaje=f"Cliente {id} no encontrado",
            codigo="CUSTOMER_NOT_FOUND",
            status_code=404
        )
    return cliente
\end{lstlisting}

% ─── SECCIÓN 4: ESTRUCTURA DE PROYECTO ───────────────────────────────────────
\section{Estructura de Proyecto Real}

\begin{lstlisting}[style=python,numbers=none]
# Estructura tipica de microservicio FastAPI en Toku
toku-invoice-service/
    src/
        main.py           # FastAPI app, startup, shutdown
        config.py         # Settings desde variables de entorno
        database.py       # Pool de conexiones asyncpg
        dependencies.py   # Deps compartidas (auth, db)
        errors.py         # Exception handlers
        routers/
            invoices.py   # POST /invoices, GET /invoices, etc.
            customers.py
            webhooks.py
        models/
            invoice.py    # Pydantic schemas
            customer.py
        services/
            invoice_service.py    # Logica de negocio
            payment_service.py
        tasks/
            payment_tasks.py      # Background tasks
    tests/
        test_invoices.py
        test_webhooks.py
    Dockerfile
    requirements.txt
\end{lstlisting}

\begin{lstlisting}[style=python]
# main.py - estructura basica
from fastapi import FastAPI
from contextlib import asynccontextmanager
from database import init_pool, close_pool

@asynccontextmanager
async def lifespan(app: FastAPI):
    # Startup
    await init_pool()
    yield
    # Shutdown
    await close_pool()

app = FastAPI(
    title="Toku Invoice Service",
    version="1.0.0",
    lifespan=lifespan
)

from routers import invoices, customers, webhooks
app.include_router(invoices.router, prefix="/api/v1")
app.include_router(customers.router, prefix="/api/v1")
app.include_router(webhooks.router)
\end{lstlisting}

% ─── SECCIÓN 5: TESTING ──────────────────────────────────────────────────────
\section{Testing de APIs FastAPI}

\begin{lstlisting}[style=python]
# tests/test_invoices.py
import pytest
from httpx import AsyncClient, ASGITransport
from unittest.mock import AsyncMock, patch
from main import app

@pytest.fixture
async def client():
    async with AsyncClient(
        transport=ASGITransport(app=app),
        base_url="http://test"
    ) as c:
        yield c

@pytest.fixture
def headers():
    return {"Authorization": "Bearer test_api_key"}

@pytest.mark.asyncio
async def test_crear_invoice_exitoso(client, headers):
    payload = {
        "customer_id": "cust_123",
        "monto": 50000,
        "moneda": "CLP",
        "fecha_vencimiento": "2027-12-31"
    }
    with patch("routers.invoices.get_current_org", return_value={"id": "org_1"}):
        with patch("routers.invoices.db") as mock_db:
            mock_db.fetchrow = AsyncMock(return_value={
                "id": "inv_new", "customer_id": "cust_123",
                "monto": 50000, "moneda": "CLP",
                "estado": "pending", "fecha_vencimiento": "2027-12-31",
                "created_at": "2026-03-17T00:00:00"
            })

            response = await client.post("/api/v1/invoices", json=payload, headers=headers)

    assert response.status_code == 201
    data = response.json()
    assert data["estado"] == "pending"
    assert data["monto"] == 50000

@pytest.mark.asyncio
async def test_crear_invoice_monto_invalido(client, headers):
    payload = {"customer_id": "c1", "monto": -100, "fecha_vencimiento": "2027-12-31"}
    response = await client.post("/api/v1/invoices", json=payload, headers=headers)
    assert response.status_code == 422
\end{lstlisting}

\section{Cheatsheet FastAPI}

\begin{multicols}{2}
\textbf{\textcolor{fapigreen}{Rutas comunes}}
\begin{lstlisting}[style=python,numbers=none]
@router.get("/items/{id}", response_model=R)
@router.post("/items", status_code=201)
@router.put("/items/{id}")
@router.delete("/items/{id}", status_code=204)
\end{lstlisting}

\textbf{\textcolor{fapigreen}{Pydantic v2}}
\begin{lstlisting}[style=python,numbers=none]
class M(BaseModel):
    x: int = Field(gt=0)
    y: Optional[str] = None
    model_config = {"from_attributes": True}
\end{lstlisting}

\textbf{\textcolor{fapigreen}{HTTPException}}
\begin{lstlisting}[style=python,numbers=none]
raise HTTPException(
    status_code=404,
    detail="No encontrado"
)
\end{lstlisting}

\columnbreak

\textbf{\textcolor{fapigreen}{Dependencies}}
\begin{lstlisting}[style=python,numbers=none]
async def dep():
    yield valor   # context manager
# uso:
def ruta(x = Depends(dep)):
    ...
\end{lstlisting}

\textbf{\textcolor{fapigreen}{Background tasks}}
\begin{lstlisting}[style=python,numbers=none]
def ruta(bg: BackgroundTasks):
    bg.add_task(fn, arg1, arg2)
    return {"status": "ok"}
\end{lstlisting}

\textbf{\textcolor{fapigreen}{Lifespan (startup/shutdown)}}
\begin{lstlisting}[style=python,numbers=none]
@asynccontextmanager
async def lifespan(app):
    # startup
    yield
    # shutdown
app = FastAPI(lifespan=lifespan)
\end{lstlisting}
\end{multicols}

\end{document}
